% ==============================================================================
% reports/chapters/01_gioi_thieu.tex
% CHƯƠNG 1: GIỚI THIỆU VÀ ĐẶT VẤN ĐỀ
% ==============================================================================

\chapter{Giới Thiệu và Đặt Vấn Đề}
\label{chap:gioi_thieu}

\section{Bối cảnh an ninh mạng và nguy cơ tấn công brute force}

Trong hạ tầng công nghệ thông tin hiện đại, an toàn an ninh thông tin là yêu cầu thiết yếu đối với các tổ chức. Trong số các hình thức tấn công mạng phổ biến, tấn công dò quét mật khẩu (Brute Force) nhắm vào các giao thức quản trị và truyền tệp từ xa như SSH (Cổng 22) và FTP (Cổng 21) vẫn là mối đe dọa thường trực. Bằng cách tự động hóa các lượt thử thông tin xác thực, kẻ tấn công có thể chiếm đoạt quyền điều khiển máy chủ nội bộ và tạo bàn đạp mở rộng tấn công trong mạng.

Các hệ thống phát hiện xâm nhập mạng (NIDS) truyền thống dựa trên luật/chữ ký như Snort \cite{roesch1999snort} hay Suricata \cite{suricata_oisf} hoạt động dựa trên đối khớp chuỗi nhị phân (chữ ký tĩnh) hoặc các quy tắc ngưỡng tĩnh. Mặc dù có ưu điểm về tốc độ đối với các mẫu đã biết, cách tiếp cận này gặp khó khăn trước các biến thể tấn công mới, chi phí duy trì luật thủ công lớn và đặc biệt bị hạn chế trước các giao thức mã hóa tầng ứng dụng như SSH, nơi nội dung tải trọng dữ liệu được bảo vệ bằng mật mã.

\section{Học máy trong phát hiện xâm nhập dựa trên luồng mạng}

Để giải quyết bài toán trên, hướng tiếp cận dựa trên học máy kết hợp phân tích luồng mạng được nghiên cứu rộng rãi. Thay vì phân tích gói tin chuyên sâu (DPI) tốn kém tài nguyên và bị cản trở bởi mã hóa, phương pháp luồng mạng chỉ trích xuất các đặc trưng thống kê tổng hợp của cả phiên truyền thông như: tổng thời lượng phiên, số lượng gói tin hai chiều, thống kê kích thước gói tin, thời gian trễ giữa các gói tin liên tiếp (IAT) và các cờ điều khiển TCP.

Các thuật toán học máy dạng cây (Decision Tree, Random Forest, XGBoost) thường được lựa chọn trong bài toán này nhờ tốc độ huấn luyện/suy luận nhanh, khả năng xử lý tốt dữ liệu dạng bảng và khả năng trích xuất các quy tắc phân loại rõ ràng.

\section{Vấn đề đánh giá: Độ lệch đánh giá và yếu tố thời gian}

Mặc dù nhiều công bố học thuật báo cáo điểm số F1 rất cao ($> 99\%$) cho các mô hình NIDS, việc áp dụng thực tế vẫn gặp nhiều thách thức. Sommer và Paxson \cite{sommer2010outside} đã nhấn mạnh sự khác biệt căn bản của an ninh mạng: môi trường mạng liên tục biến động và kẻ tấn công chủ động thay đổi hành vi.

Gần đây, Arp và cộng sự \cite{arp2022dos} đã chỉ ra nhiều cạm bẫy thực nghiệm phổ biến trong học máy cho an ninh mạng, nổi bật là \textbf{độ lệch đánh giá quá lạc quan} khi áp dụng chiến lược \textbf{phân chia dữ liệu ngẫu nhiên}.

Khi áp dụng phân chia ngẫu nhiên trên tập dữ liệu mạng như CICIDS2017 \cite{sharafaldin2018generating}:
\begin{itemize}
    \item Các flow thuộc cùng một chiến dịch tấn công (phát sinh từ cùng một công cụ, cùng địa chỉ IP nguồn, cùng cấu hình thử nghiệm) có thể bị phân bổ đồng thời vào cả tập Huấn luyện, Kiểm định và Kiểm thử.
    \item Sự tương quan mạnh nội tại giữa các mẫu trong cùng chiến dịch có thể khiến giả thiết các quan sát độc lập và phân phối đồng nhất ($i.i.d.$) không còn được đảm bảo chặt chẽ.
    \item Do đó, mô hình có thể đạt điểm số cao nhờ việc nhận dạng lại dấu hiệu thống kê của chiến dịch đã thấy (mang tính chất nội suy) thay vì năng lực tổng quát hóa sang chiến dịch mới.
\end{itemize}

\section{Mục tiêu và câu hỏi nghiên cứu}

Bài tập lớn này xây dựng một quy trình thực nghiệm trên tập dữ liệu CICIDS2017 Tuesday, so sánh các giao thức phân tách theo thời gian và ngẫu nhiên, đồng thời khảo sát cách các mô hình học máy dạng cây xử lý biến thể SSH không có trong tập fit của benchmark chuẩn. Do SSH xuất hiện trong Validation và nhãn của các mẫu này tham gia lựa chọn cấu hình, thiết lập không phải quy trình zero-shot khép kín từ đầu đến cuối.

Đề tài giải quyết 4 câu hỏi nghiên cứu (RQ) sau:
\begin{itemize}
    \item[\textbf{RQ1:}] \textit{(Đánh giá trên subtype SSH vắng mặt trong tập fit theo thời gian)} Khi mô hình dạng cây được huấn luyện trên dữ liệu Train chỉ chứa \FTP, khả năng phát hiện \SSH\ xuất hiện ở giai đoạn thời gian sau thay đổi như thế nào?
    \item[\textbf{RQ2:}] \textit{(So sánh các giao thức phân tách)} Kết quả của Random Split và Time-based Split khác nhau như thế nào khi đồng thời xét đến chiến lược chia, thành phần subtype và phân bố tập kiểm thử?
    \item[\textbf{RQ3:}] \textit{(Vai trò của cổng dịch vụ và các đặc trưng hành vi)} Cổng dịch vụ đích (\texttt{Destination Port}) đóng vai trò như thế nào trong dự đoán của mô hình? Khi loại bỏ cổng mạng, các đặc trưng thống kê hành vi gói tin có đủ để nhận diện lưu lượng tấn công hay không?
    \item[\textbf{RQ4:}] \textit{(So sánh hai pipeline XGBoost)} Hai pipeline XGBoost được ghi nhận khác nhau như thế nào về dữ liệu fit, đặc trưng, tiền xử lý và cấu hình; những khác biệt này giới hạn ra sao khả năng quy kết chênh lệch kết quả cho chính sách refit?
\end{itemize}

\section{Đóng góp thực nghiệm của bài tập lớn}

Các kết quả thực nghiệm chính của bài tập lớn bao gồm:
\begin{enumerate}
    \item \textbf{Xây dựng quy trình thực nghiệm theo thời gian:} Thực hiện tiền xử lý và tái dựng timestamp 24 giờ từ CSV; áp dụng Purge ($q=60\text{s}$) và Embargo ($g=120\text{s}$), loại bỏ \DataPurgedRows\ dòng tại vùng ranh giới để giảm nguy cơ flow vượt mốc cắt. Benchmark chuẩn fit mô hình trên Train-only (kịch bản FTP $\to$ SSH); Validation có \DataValidationSSH\ mẫu SSH và nhãn của chúng tham gia lựa chọn cấu hình, nhưng không tham gia fit trọng số mô hình benchmark.
    \item \textbf{Đánh giá thực nghiệm 3 mô hình học máy dạng cây:} Đối chuẩn Decision Tree, Random Forest và XGBoost qua 12 tổ hợp thực nghiệm (3 mô hình $\times$ 2 chiến lược chia $\times$ 2 kịch bản cổng). Ghi nhận trung thực sự suy giảm hiệu năng của mô hình trên tập Test theo thời gian: trong kịch bản chính có cổng mạng, Random Forest đạt $\FOne = \RFTestFOne$, nhận diện được \RFTestSSHRecall\ tấn công SSH (\RFTestTP/\DataTestSSH\ flow).
    \item \textbf{Phân tích cơ chế suy luận qua SHAP và cấu trúc cây:} Sử dụng SHAP TreeExplainer và phân tích nhánh rẽ của cây quyết định để làm rõ việc mô hình phụ thuộc mạnh vào cổng dịch vụ (có giá trị SHAP trung bình lớn nhất \SHAPTopOneVal, gấp \SHAPRatioTopOneTwo\ lần đặc trưng thứ hai).
    \item \textbf{Đối chiếu hai pipeline XGBoost:} Phân tích khác biệt về ranh giới fit, tập đặc trưng, tiền xử lý và cấu hình giữa benchmark chuẩn và lần chạy tuning độc lập. Lần chạy tuning dùng \texttt{refit=True} theo đúng thiết kế của scikit-learn và fit mô hình cuối cùng trên Train + Validation, trong đó có \DataValidationSSH\ mẫu SSH; do các pipeline khác nhau ở nhiều yếu tố, chênh lệch Test F1 không được quy riêng cho refit.
    \item \textbf{Tổ chức mã nguồn và số liệu có tính tái lập cao:} Toàn bộ số liệu trong báo cáo được truy xuất tự động từ các artifact đã đóng băng, đi kèm mã băm SHA-256 xác minh tệp nguồn, với khả năng tái lập được kiểm soát thông qua cấu hình môi trường đã ghi nhận và kiểm chứng (xem Phụ lục \ref{app:reproducibility}).
\end{enumerate}

\section{Cấu trúc của báo cáo}

Báo cáo được tổ chức thành 7 chương:
\begin{itemize}
    \item \textbf{Chương 1:} Giới thiệu bối cảnh an ninh mạng, vấn đề khoa học, mục tiêu và 4 câu hỏi nghiên cứu.
    \item \textbf{Chương 2:} Cơ sở lý thuyết về luồng mạng, cơ chế tấn công Brute Force, các thuật toán học máy dạng cây và hệ thống chỉ số đánh giá mất cân bằng lớp.
    \item \textbf{Chương 3:} Chi tiết tập dữ liệu CICIDS2017 Tuesday, mã băm SHA-256, quy tắc phục hồi thời gian và quy trình làm sạch dữ liệu.
    \item \textbf{Chương 4:} Phương pháp đề xuất và kiến trúc hệ thống, chính sách đặc trưng, kỹ thuật Purge/Embargo và nguyên tắc đóng băng cấu hình.
    \item \textbf{Chương 5:} Thiết kế thực nghiệm, không gian siêu tham số và tiêu chí lựa chọn cấu hình mô hình trên tập kiểm định.
    \item \textbf{Chương 6:} Kết quả và phân tích thực nghiệm qua 12 tổ hợp, giải thích SHAP, cấu trúc rẽ nhánh cây và nghiên cứu điển hình về cờ refit.
    \item \textbf{Chương 7:} Kết luận, trả lời 4 câu hỏi nghiên cứu, đánh giá tính hợp lệ, hướng phát triển tương lai và bảng phân công nhiệm vụ thành viên.
\end{itemize}
